\section{Konzept und Pipeline-Idee}

\subsection{Leitprinzip und Datenfluss}

Die Pipeline trennt die Bilddomäne, die semantische Domäne und die
Geometriedomäne. Der Ablauf ist:

\begin{enumerate}
  \item Aus dem Rohvideo werden Arbeitsframes erzeugt.
  \item SAM~3.1 propagiert den Objektprompt und speichert hierarchische Masken.
  \item COLMAP schätzt aus unmaskierten Frames ein ursprüngliches
        Kameramodell und eine sparse Punktwolke.
  \item Bei verzeichnungsbehafteten Kameramodellen wird eine ideale PINHOLE-
        Szene erzeugt; die Masken werden mit derselben Abbildung gewarpt.
  \item STS trainiert die objektspezifische Gaussian-Repräsentation.
  \item Die Hochopazitätskopie der Objekt-Gaussians wird über die gewählte
        Meshroute in ein Coarse-Mesh überführt.
  \item DGtal extrahiert daraus eine lokale Centerline; anschließend können
        B-Spline und GeoJSON erzeugt werden.
\end{enumerate}

Abbildung~\ref{fig:pa-pipeline-flow} fasst diesen Datenfluss zusammen. Sie ist
aus dem Flowchart des Exposés abgeleitet, wurde für den aktuellen
Projektarbeitsstand jedoch angepasst: Die Web-UI und der nicht systematisch
untersuchte SAHI-Fallback sind nicht Bestandteil des Diagramms. Stattdessen
werden die beiden aktuellen Meshrouten A und \texttt{sugar\_coarse} explizit
getrennt dargestellt. Der CloudCompare-Schritt ist als optionaler
Georeferenzierungszweig enthalten und nicht als validierter PA-
Genauigkeitsnachweis zu lesen.

\begin{figure}[H]
\centering
\resizebox{\textwidth}{!}{%
\begin{tikzpicture}[
      node distance=1.35cm and 2.0cm,
      box/.style={rectangle, draw=blue!65, fill=blue!5, very thick,
            minimum width=4.4cm, minimum height=0.9cm, align=center,
            rounded corners, text width=4.2cm},
      route/.style={rectangle, draw=purple!70, fill=purple!7, very thick,
            minimum width=4.6cm, minimum height=0.9cm, align=center,
            rounded corners, text width=4.4cm},
      decision/.style={diamond, draw=red!65, fill=red!7, very thick,
            text width=2.7cm, align=center, aspect=2.0},
      arrow/.style={-{Stealth[scale=1.05]}, thick},
      container/.style={rectangle, draw, dashed, rounded corners, inner sep=0.35cm}
]
      \node[box] (input) {Rohvideo\\Arbeitsframes\\Prompt \texttt{pipe}};
      \node[box, below=of input] (sam) {\textbf{Container A: SAM 3.1}\\Temporale Segmentierung\\default / middle / small};
      \node[box, below=of sam] (colmap) {\textbf{Container B: COLMAP}\\Unmaskiertes SfM\\Originales Kameramodell und Sparse-Modell};
      \node[box, below left=1.55cm and 2.4cm of colmap] (ideal) {Idealbilddomäne\\\texttt{image\_undistorter}\\Maskenwarp mit identischer Abbildung};
      \node[box, below right=1.55cm and 2.4cm of colmap] (gcp) {Originalmodell\\GCP-/CloudCompare-Zweig\\Matrix optional und PA-seitig nicht validiert};
      \node[box, below=1.55cm of ideal] (sts) {\textbf{Container C: STS}\\Objekt-ID-Zuordnung\\Objektbezogene Gaussian-Repräsentation};
      \node[decision, below=1.55cm of sts] (meshroute) {Meshroute};
      \node[route, below left=1.55cm and 1.7cm of meshroute] (original) {\textbf{Route A: Original-GS}\\Surface-Sampling, Poisson, Decimation\\kein SuGaR-Coarse/Refinement/UV};
      \node[route, below right=1.55cm and 1.7cm of meshroute] (coarse) {\textbf{SuGaR-Coarse-Vergleich}\\mask-aware Fork, Coarse, optional Refinement/UV};
      \node[box, below=2.0cm of meshroute] (post) {\textbf{Container E: DGtal/GDAL}\\lokales Mesh $\rightarrow$ Centerline\\B-Spline $\rightarrow$ lokale CSV/GeoJSON};
      \node[box, below right=1.8cm and 1.8cm of post] (geo) {Optionaler Exportzweig\\4x4-Transformation + Anchor\\UTM/GIS nur mit validierter Matrix};

      \begin{pgfonlayer}{background}
            \node[container, fill=blue!8, draw=blue!75, fit=(input) (sam)] {};
            \node[container, fill=orange!8, draw=orange!75, fit=(colmap)] {};
            \node[container, fill=cyan!8, draw=cyan!75, fit=(ideal) (sts)] {};
            \node[container, fill=red!8, draw=red!75, fit=(gcp)] {};
            \node[container, fill=purple!8, draw=purple!75, fit=(meshroute) (original) (coarse)] {};
            \node[container, fill=teal!8, draw=teal!75, fit=(post) (geo)] {};
      \end{pgfonlayer}

      \draw[arrow] (input) -- (sam);
      \draw[arrow] (sam) -- (colmap);
      \draw[arrow] (colmap) -- (ideal);
      \draw[arrow] (colmap) -- (gcp);
      \draw[arrow] (ideal) -- (sts);
      \draw[arrow] (sts) -- (meshroute);
      \draw[arrow] (meshroute) -- node[above left] {A} (original);
      \draw[arrow] (meshroute) -- node[above right] {Vergleich} (coarse);
      \draw[arrow] (original) |- (post);
      \draw[arrow] (coarse) |- (post);
      \draw[arrow] (post) -- (geo);
      \draw[arrow] (gcp) |- (geo);
\end{tikzpicture}%
}
\caption{Aktualisierter Datenfluss der fünf Container mit Original-GS-
Produktionsroute, SuGaR-Coarse-Vergleich und optionalem
Georeferenzierungszweig.}
\label{fig:pa-pipeline-flow}
\end{figure}

Das Konzept verfolgt zwei gleichrangige Ziele. Erstens soll ein geführter
Einzellauf von der Eingabe bis zum Endprodukt ohne manuelle Neuverkabelung der
Werkzeuge möglich sein. Zweitens müssen teure Stufen kontrolliert wiederholt
und als Matrix ausgeführt werden können. Jeder Übergang besitzt deshalb
prüfbare Eingaben und persistente Ausgaben.

\subsection{Domänentrennung als zentraler Datenvertrag}

Die methodische Kernidee ist die Domänentrennung: GCP und SfM bleiben im
Originalraum, STS und SuGaR arbeiten in einer idealisierten, konsistenten
Bilddomäne. Dadurch wird vermieden, dass eine Bildverzeichnung nur in den
Kameraparametern, nicht aber in den Masken berücksichtigt wird.

Bei \texttt{SIMPLE\_RADIAL} und \texttt{OPENCV} wird zuerst ein originales
COLMAP-Modell auf den Rohframes geschätzt. Der Undistorter erzeugt daraus eine
separate ideale Szene. Die Masken werden mit der identischen
Kameramodellabbildung transformiert. Bei \texttt{PINHOLE} wird dieselbe
Schnittstelle als Pass-through verwendet. Somit sehen STS, Renderer,
SuGaR-Route und Metrikberechnung stets zueinander passende Pixel- und
Kamerakoordinaten.

Diese vollständige Domänenkette ist im automatisierten Matrixpfad sowie im
restaurierten Smoke-/Replaypfad nachgewiesen. Der historische Inline-Vollauf
und der allgemeine \texttt{--from}-Dispatcher erzeugen zwar die ideale
COLMAP-Szene, rufen den separaten Maskenwarp derzeit jedoch nicht selbst auf.
Bis zur Vereinheitlichung des Codes beziehen sich Aussagen über die vollständig
gewarpten Bild-/Maskendomäne daher ausdrücklich auf Matrix- und Replayläufe.

\subsection{Semantische Objektkette}

Die SAM-Masken erfüllen mehrere getrennte Funktionen. Sie definieren die
Objektzuordnung in STS, begrenzen die Auswahl der Ziel-Gaussians und stellen die
Objektregion für die Bildmetriken bereit. In der SuGaR-Coarse-Vergleichsroute
begrenzen sie zusätzlich die photometrische Supervision. Diese mehrfache
Nutzung darf nicht mit einer geometrischen Garantie verwechselt werden: Eine
binäre Maske beschreibt Bildpixel, während ein räumlich ausgedehnter Gaussian
und das aus Samples rekonstruierte Mesh zusätzliche Unsicherheiten besitzen.

\subsection{Produktions- und Vergleichsroute}

Für die Projektarbeit wird Variante A als Produktionsroute verwendet. SuGaR-
Coarse und Refinement bleiben als Vergleichsroute erhalten, sind aber nicht
notwendig, um die direkte STS-Gaussian-Geometrie zu vermessen. Das Mesh ist ein
abgeleitetes Produkt; die Gaussian-Wolke bleibt die primäre geometrische
Hypothese.

Route A führt keinen neuen SuGaR-Splat ein. Sie extrahiert die Oberfläche aus
der vorbereiteten Original-STS-Gaussian-Wolke. Die Vergleichsroute optimiert
dieselbe Ausgangslage zunächst mit dem lokalen maskenbewussten SuGaR-Fork und
kann anschließend ein gebundenes Refinement und UV-Baking ausführen. Die
Trennung ermöglicht, Änderungen der Gaussian-Geometrie von Änderungen der
eigentlichen Mesh-Extraktion zu unterscheiden.

\subsection{Ausführungsformen}

Der interaktive Vollauf führt durch Eingabeauswahl und Parameterentscheidungen.
Der Autopilot übernimmt validierte Defaults und überspringt den manuellen
CloudCompare-Breakpoint, wenn keine echte Transformationsmatrix bereitsteht.
Der Replay-Pfad startet an einer festgelegten Schrittgrenze mit vorhandenen,
konsistenten Artefakten. Er setzt kein abgebrochenes Training innerhalb einer
Iteration fort. Der Matrixrunner führt deklarierte Kombinationen
stdin-unabhängig nacheinander aus, archiviert jeden Lauf und setzt den
Live-Arbeitsbereich zurück. Diese Modi nutzen dieselben fachlichen
Pipelinekomponenten, erfüllen aber verschiedene Aufgaben der Benutzerführung
und Reproduzierbarkeit.

Die Projektarbeit untersucht die CLI-zentrierte Ausführung. Eine vorhandene
Web-UI und ein tatsächliches Serverdeployment bleiben außerhalb des Scopes.
Die Containerisierung und die nichtinteraktive Matrix zeigen eine für einen
späteren Serverbetrieb geeignete Architektur, aber noch keinen validierten
Produktiv-Rollout.

\subsection{Lokaler Endpunkt und vorbereitete Georeferenzierung}

Die Georeferenzierung ist als nachgelagerter Datenvertrag vorgesehen. Die
Centerline wird zunächst lokal erzeugt. Erst eine validierte Transformation aus
GCP-Picking und ein separater Anchor überführen sie in ein UTM-System. Ein
Fallback mit reiner Translation dient nur der technischen Pipelineprüfung und
nicht als Genauigkeitsnachweis.
